\chapter{Introduction}
\label{ch:intro}

Heavy ion collisions are a unique laboratory setting developed over the last 50 years to study in a controlled manner what nuclear matter is like in extreme conditions. 
The construction and operation of heavy ion collision experiments developed out of a desire to understand questions across a wide range of physics: from understanding more about the early universe, to better understanding the nuclear strong force and mechanism for astrophysical showers of particles produced in cosmic ray interactions. 
For questions posed in the mid 20th century, including "what was nuclear matter in the early universe like?" and "what is the nature of the strong nuclear force in a collective system?", experimental results from heavy ion collisions have been able to provide answers and open up additional discoveries. 

At everyday temperatures and densities of our world, nuclear matter is organized into hadrons, such as protons and neutrons. 
However, these hadrons are not fundamental particles. They are instead made up of quarks and gluons, which are fundamental particles that interact via the nuclear strong force. 
The strong nuclear force is one of the four fundamental forces within the Standard Model of particle physics along with the weak, electromagnetic and gravitational forces.
However, the governing theory of the strong force, quantum chromodynamics (QCD), is difficult to study as it is not directly accessible in ordinary matter.
This is because only quarks and gluons carry color charge and therefore interact via the strong force; the hadrons themselves are color neutral, making it impossible to probe the strong force directly in ordinary matter. 
Therefore, studying these fundamental constituents of nuclear matter required developing experiments capable of reaching extreme energies where these quarks and gluons can be liberated from their parent hadrons.

Heavy ion collisions have been studied at a variety of fixed target experiments and hadron colliders, including the Alternating Gradient Synchrotron (AGS) at BNL in the 1980s, the Super Proton Synchrotron (SPS) at CERN in the 1990s, the Relativistic Heavy Ion Collider (RHIC) at BNL from the 2000s onwards, and the Large Hadron Collider (LHC) at CERN from the 2010s onwards.
These collisions have used a variety of different ion species, from protons to light ions, like oxygen and nitrogen, and heavy ions, like gold and lead. 
In these collisions, these ions are accelerated in rings of magnets that are miles long to nearly the speed of light and then collided together or into a fixed target. 
These collisions are very short, on the order of $10^{-23}$ seconds, but they can produce very high energy densities in the collision region, with temperatures reaching 250,000 times hotter than the sun in RHIC collisions.
Here, a key motivation for using heavy ions in these collisions is their larger size which leads to very high energy densities in the collision region and provides a finite system size for the  study of collective nuclear matter effects rather than just individual nucleon-nucleon interactions.

\begin{figure}
\centering
\includegraphics[width=\textwidth]{Figures/history-of-the-universe-2015.pdf}
\caption{A timeline of the history of the universe. Time scales around the "NUCLEONS FORM" region are probed via heavy ion collisions at RHIC and the LHC are shown in the top and bottom labels (t = $10^{-4}$ s and E = $10^{-1}$ GeV). Image from LBNL.}
\end{figure}

At the very high energies reached in these collisions, the quarks and gluons are then liberated from their parent hadrons and form a state of matter called the quark gluon plasma (QGP). 
The QGP is believed to have existed in the early universe, a few microseconds after the Big Bang.
At this time, the universe had cooled enough for fundamental particles to acquire mass via the electroweak phase transition, while the strong phase transition, during which quarks and gluons became bound into hadrons and created the nuclear matter in our universe today, was still ongoing.
Therefore, heavy ion collisions provide a unique opportunity to recreate the conditions of this stage of the early universe and study this phase of matter in a laboratory setting.

\begin{figure}
\centering
\includegraphics[width=0.7\textwidth]{Figures/phase_transition.jpg}
\caption{Phase diagram of nuclear matter with axes of temperature and baryon density. The regions of nuclear matter that are probed by heavy ion collisions at the LHC (green) and RHIC (blue) are shown.. Image from BNL}
\end{figure}

Once it became clear that heavy ion collisions could probe this new state of matter, research questions further evolved from "are we creating the QGP?" to "what does this QGP actually look like?"
Surprising early results from heavy ion collisions at RHIC showed that the QGP is not a weakly interacting gas of quarks and gluons as originally expected, but instead behaves as a strongly interacting fluid with very low specific viscosity.
Additional results showed that high energy particles produced in the initial state of the collision were modified by interactions with the QGP in a process called “jet quenching”. 
Present questions in the field continue to focus on understanding and quantifying the properties of the QGP, such as its transport properties, as well as better understanding the initial state of the colliding nuclei and how that affects our understanding of the evolution of the collision and the properties of the QGP.

This thesis presents two measurements from the sPHENIX experiment at RHIC that contribute to our understanding of the questions just presented about the quark gluon plasma and the early universe. The first measurement on the transverse energy density in Au+Au collisions at $\sqrt{s_{NN}}$ = 200 GeV is a direct measurement of the energy density of the quark gluon plasma. This measurement has historically been used as evidence for QGP formation, experimentally proving that the energy density is well above the strong phase transition energy density predicted to yield QGP. Here, this measurement is used to validate the sPHENIX calorimeter performance, reconstruction and calibration as the first calorimeter physics measurement using the sPHENIX detector. This measurement also provides a more precise estimate of the energy density than previous measurements and is directly compared to theoretical predictions. 
The second measurement on the underlying event production in p+p collisions at $\sqrt{s}$ = 200 GeV further investigates the initial state of heavy ion and proton-proton collisions, a part of the collision model that is still not well understood. This measurement provides a probe of correlations between hard and soft processes in p+p collisions and can help distinguish between best models of the initial state in these collisions.

This thesis is organized as follows: Chapter 2 provides theoretical background on QCD, the QGP, and the modeling of heavy ion collisions. Chapter 3 describes the experimental setup of the sPHENIX experiment at RHIC. Chapter 4 outlines commissioning, reconstruction and calibration of the sPHENIX calorimeters. Chapter 5 presents the measurement of the transverse energy density in Au+Au collisions at $\sqrt{s_{NN}}$ = 200 GeV. Chapter 6 presents the measurement of the underlying event production in p+p collisions at $\sqrt{s}$ = 200 GeV. Finally, Chapter 7 summarizes and discusses these results in the context of the field of heavy ions.